%

\index{zadanie!Monge'a}
\pol{\subsection{Problem Monge'a}}
\ita{\subsection{Problema di Monge}}

\begin{problem}[\pol{Monge'a?}\ita{di Monge?}]
    \pol{Okrąg, który przecina trzy okręgi $\Gamma_1$, $\Gamma_2$, $\Gamma_3$ pod kątem prostym.}
    \ita{Costruire un cerchio che intersechi ortogonalmente tre cerchi $\Gamma_1$, $\Gamma_2$, $\Gamma_3$.}
\end{problem}

% https://mathworld.wolfram.com/MongesProblem.html
\pol{Dla każdej pary okręgów $\Gamma_i$, $\Gamma_j$ znajdujemy oś potęgową; jeśli trzy osie przecinają się w punkcie $O$ leżącym na zewnątrz okręgów $\Gamma_i$, to jest to środek szukanego okręgu.}
\ita{Per ogni coppia di cerchi $\Gamma_i$, $\Gamma_j$ si determina l'asse radicale; se i tre assi si intersecano in un punto $O$ situato all'esterno dei cerchi $\Gamma_i$, allora esso è il centro del cerchio cercato.}
\pol{Promieniem szukanego okręgu jest odcinek styczny do $\Gamma_i$ oraz przechodzący przez $O$.}
\ita{Il raggio del cerchio cercato è dato dal segmento tangente a $\Gamma_i$ passante per $O$.}
\pol{Jeśli jednak środek okręgu leży wewnątrz któregoś okręgu albo osie nie przecinają się, to problem nie ma rozwiąania.}
\ita{Se invece il centro del cerchio giace all'interno di uno dei cerchi oppure gli assi non si intersecano, allora il problema non ha soluzione.}

%